\bookchapter{Matrices}\label{ch:matrices}

\section*{Explanation}
A \emph{matrix} is a rectangular table of numbers. A matrix with $m$ rows and $n$ columns is an $m\times n$ matrix, and the entry in row $i$ and column $j$ is written $a_{ij}$. For example
\[
A=\begin{pmatrix}1&2\\3&4\\5&6\end{pmatrix}
\]
is a $3\times2$ matrix with $a_{21}=3$. A vector is simply a matrix with one column.

\section*{Matrix--vector product}
To multiply an $m\times n$ matrix $A$ with a vector $\vv x$ of length $n$, we compute the dot product of \emph{every row} of $A$ with $\vv x$. The result is a vector of length $m$:
\[
\begin{pmatrix}1&2\\3&4\\5&6\end{pmatrix}\begin{pmatrix}2\\1\end{pmatrix}
=\begin{pmatrix}1\cdot2+2\cdot1\\3\cdot2+4\cdot1\\5\cdot2+6\cdot1\end{pmatrix}
=\begin{pmatrix}4\\10\\16\end{pmatrix}.
\]
\begin{keyidea}[Shapes must fit]
$(m\times n)\cdot(n\times1)=(m\times1)$. The number of columns of the matrix must equal the length of the vector.
\end{keyidea}

\section*{Matrix--matrix product}
The entry in row $i$ and column $j$ of $AB$ is the dot product of \emph{row $i$ of $A$} with \emph{column $j$ of $B$}. The shapes must fit: $(m\times n)\cdot(n\times p)=(m\times p)$.
\[
\begin{pmatrix}\textcolor{red}{1}&\textcolor{red}{2}\\3&4\end{pmatrix}
\begin{pmatrix}\textcolor{red}{0}&1\\\textcolor{red}{1}&0\end{pmatrix}
=\begin{pmatrix}\textcolor{red}{2}&1\\4&3\end{pmatrix}
\qquad(\text{red: }1\cdot0+2\cdot1=2).
\]
The order matters: for the same two matrices,
\[
AB=\begin{pmatrix}2&1\\4&3\end{pmatrix}\quad\text{(swaps the columns of }A),\qquad
BA=\begin{pmatrix}3&4\\1&2\end{pmatrix}\quad\text{(swaps the rows of }A),
\]
so in general $AB\neq BA$.

\begin{example}
\[
\begin{pmatrix}1&0&2\\-1&3&1\end{pmatrix}\begin{pmatrix}3&1\\2&1\\1&0\end{pmatrix}
=\begin{pmatrix}5&1\\4&2\end{pmatrix}.
\]
Here $(2\times3)\cdot(3\times2)=(2\times2)$. For instance, the entry $4$ in row 2, column 1 is $(-1)\cdot3+3\cdot2+1\cdot1=4$.
\end{example}

\section*{Transpose and identity}
The \emph{transpose} $A^{\top}$ swaps rows and columns: $(A^\top)_{ij}=a_{ji}$.
\[
\begin{pmatrix}1&2&3\\4&5&6\end{pmatrix}^{\top}=\begin{pmatrix}1&4\\2&5\\3&6\end{pmatrix},\qquad
(AB)^\top=B^\top A^\top.
\]
The \emph{identity matrix} $I$ has ones on the diagonal and zeros elsewhere. It does nothing: $IA=AI=A$ and $I\vv x=\vv x$. For $2\times2$: $I=\begin{pmatrix}1&0\\0&1\end{pmatrix}$.

\section*{A matrix transforms the plane}
A $2\times2$ matrix $A$ turns every point $\vv x$ into the point $A\vv x$. The columns of $A$ are the images of the basis vectors $\vv e_1=(1,0)$ and $\vv e_2=(0,1)$. So we can \emph{see} what a matrix does by looking at what happens to the unit square.

\begin{center}
\begin{tikzpicture}[>=Stealth, scale=0.85]
  \begin{scope}[xshift=0cm]
    \draw[gray!40,very thin] (-0.5,-0.5) grid (2.5,2.5);
    \draw[->] (-0.5,0) -- (2.7,0); \draw[->] (0,-0.5) -- (0,2.7);
    \draw[dashed,thick] (0,0) -- (1,0) -- (1,1) -- (0,1) -- cycle;
    \filldraw[fill=blue!20,draw=blue!70!black,thick,fill opacity=0.6] (0,0) -- (2,0) -- (2,1) -- (0,1) -- cycle;
    \draw[->,very thick,red] (0,0) -- (2,0);
    \draw[->,very thick,green!50!black] (0,0) -- (0,1);
    \node at (1,-1) {$\begin{pmatrix}2&0\\0&1\end{pmatrix}$};
  \end{scope}
  \begin{scope}[xshift=4.2cm]
    \draw[gray!40,very thin] (-0.5,-0.5) grid (2.5,2.5);
    \draw[->] (-0.5,0) -- (2.7,0); \draw[->] (0,-0.5) -- (0,2.7);
    \draw[dashed,thick] (0,0) -- (1,0) -- (1,1) -- (0,1) -- cycle;
    \filldraw[fill=blue!20,draw=blue!70!black,thick,fill opacity=0.6] (0,0) -- (1,0) -- (2,1) -- (1,1) -- cycle;
    \draw[->,very thick,red] (0,0) -- (1,0);
    \draw[->,very thick,green!50!black] (0,0) -- (1,1);
    \node at (1,-1) {$\begin{pmatrix}1&1\\0&1\end{pmatrix}$};
  \end{scope}
  \begin{scope}[xshift=9.4cm]
    \draw[gray!40,very thin] (-1.5,-0.5) grid (1.5,2.5);
    \draw[->] (-1.5,0) -- (1.7,0); \draw[->] (0,-0.5) -- (0,2.7);
    \draw[dashed,thick] (0,0) -- (1,0) -- (1,1) -- (0,1) -- cycle;
    \filldraw[fill=blue!20,draw=blue!70!black,thick,fill opacity=0.6] (0,0) -- (0,1) -- (-1,1) -- (-1,0) -- cycle;
    \draw[->,very thick,red] (0,0) -- (0,1);
    \draw[->,very thick,green!50!black] (0,0) -- (-1,0);
    \node at (0,-1) {$\begin{pmatrix}0&-1\\1&0\end{pmatrix}$};
  \end{scope}
\end{tikzpicture}

\small\textbf{Figure 7.} What three matrices do to the unit square (dashed): stretching in $x$, shearing, and rotating by $90^\circ$. The red arrow is the image of $\vv e_1$ (first column), the green arrow is the image of $\vv e_2$ (second column).
\end{center}

\section*{A layer of a neural network}
A whole layer of neurons is one matrix--vector product. If the layer has $m$ neurons and receives $n$ inputs, the weights form an $m\times n$ matrix $W$ whose row $j$ holds the weights of neuron $j$:
\[
\vv z=W\vv x+\vv b,\qquad \vv h=\sig(\vv z)\ \ (\text{the activation is applied to each entry}).
\]

\begin{example}
Two inputs $\vv x=(1,0.5)$ and two neurons with
\[
W=\begin{pmatrix}0.4&0.3\\-0.2&0.6\end{pmatrix},\qquad
W\vv x=\begin{pmatrix}0.4\cdot1+0.3\cdot0.5\\-0.2\cdot1+0.6\cdot0.5\end{pmatrix}=\begin{pmatrix}0.55\\0.10\end{pmatrix}.
\]
(These are the numbers of the tiny network example.) With the sigmoid: $\vv h\approx(0.634,0.525)$.
\end{example}

To process many examples at once we put them as columns of a matrix $X$ and compute $WX$: this is why neural networks run fast on graphics cards, which are built for matrix products.

\begin{aibox}
A network with $784$ inputs and a first layer of $128$ neurons has a weight matrix of size $128\times784$, i.e.\ $784\cdot128=100\,352$ weights. Forward propagation through the network is a chain of matrix products and activation functions. A matrix multiplication costs about one multiplication per weight.
\end{aibox}

\section*{Practice problems}
\begin{enumerate}[label=\textbf{\thechapter.\arabic*}]
  \item Let $A$ be $2\times3$, $B$ be $3\times4$ and $C$ be $2\times2$. Which of $AB$, $BA$, $CA$, $AC$, $B^\top A^\top$ are defined, and what are their sizes?
  \item Compute $A\vv x$ for $A=\begin{pmatrix}2&0\\1&3\end{pmatrix}$ and $\vv x=(1,2)$.
  \item Compute $AB$ and $BA$ for $A=\begin{pmatrix}1&2\\0&1\end{pmatrix}$ and $B=\begin{pmatrix}3&0\\1&2\end{pmatrix}$.
  \item Write down the transpose of $\begin{pmatrix}1&2&3\\4&5&6\end{pmatrix}$.
  \item A layer has $W=\begin{pmatrix}1&-1\\2&0\\0&3\end{pmatrix}$, bias $\vv b=(-2,1,-4)$ and input $\vv x=(2,1)$. Compute $\vv z=W\vv x+\vv b$ and the outputs with ReLU activation.
  \item Find a $2\times2$ matrix that doubles the $x$-coordinate and keeps $y$ unchanged. Find a matrix that swaps $x$ and $y$. Apply both to $(3,4)$.
  \item Find the images of the corners $(0,0)$, $(1,0)$, $(1,1)$, $(0,1)$ of the unit square under $A=\begin{pmatrix}1&2\\0&1\end{pmatrix}$. What are the images of $\vv e_1$ and $\vv e_2$?
  \item $\star$ Check $(AB)^\top=B^\top A^\top$ for $A=\begin{pmatrix}1&2\\3&4\end{pmatrix}$ and $B=\begin{pmatrix}0&1\\1&0\end{pmatrix}$.
  \item $\star$ How many weights and biases does a layer with $784$ inputs and $128$ neurons have in total?
\end{enumerate}

% =====================================================
